\chapter{Working with formulas}
\label{ch:formulas}

Chapter~\ref{sec:language} describes the expression language.  This
chapter gives the Scheme procedures that work with its formulas:
\begin{itemize}
\item the parser and the printer, which turn a string into a formula and
  back (Section~\ref{sec:parser-printer});
\item the accessors that name the parts of a formula
  (Section~\ref{sec:parts});
\item the \texttt{wff} record that holds a checked formula
  (Section~\ref{chap:syntax-contexts});
\item the two ways of changing how a formula is printed:
  \texttt{notation!} (Section~\ref{sec:notation}) and the optional infix
  layer (Section~\ref{sec:infix-layer}).
\end{itemize}

%% =========================================================

\section{Parser and printer}
\label{sec:parser-printer}

Four entry points provide string-level access to VNB formulas:

\begin{VNBsnip}
(make-wff-from-string str)  ; string -> <wff>  (validates via make-wff)
(parse-string str)          ; string -> raw S-expression (no validation)
(wff->string w)             ; <wff>  -> string
(pp w)                      ; display wff->string then newline
\end{VNBsnip}

\noindent
\texttt{parse-string} converts a string to a raw S-expression without
validation; it is used internally by the interactive commands and is
available for low-level \texttt{cmd-*} calls.  For normal use, pass
formula strings directly to \texttt{ai}, \texttt{bc}, \texttt{inst},
\texttt{cut}, \texttt{wk}, and \texttt{vnb-create-num-theorem}.
\texttt{make-wff-from-string} is the right choice when constructing a
formula to prove (for \texttt{sp}).

\paragraph{Operator precedence (high to low).}
\begin{center}
\small
\begin{tabular}{lll}
\toprule
Level & Operators & Associativity \\
\midrule
primary & atoms, \texttt{f(a,b)}, \texttt{[a,b]}, \texttt{(expr)} & --- \\
power & \texttt{\textasciicircum} & right \\
unary & \texttt{-x} & right \\
mul & \texttt{*}, \texttt{/} & left (n-ary \texttt{*}) \\
add & \texttt{+}, \texttt{-} & left (n-ary \texttt{+}) \\
cmp & \texttt{in}, \texttt{=}, \texttt{==}, \texttt{<=} & binary \\
and & \texttt{and} & left (n-ary) \\
or & \texttt{or} & left (n-ary) \\
iff & \texttt{iff} & left (n-ary) \\
implies & \texttt{implies} & right \\
\bottomrule
\end{tabular}
\end{center}

\paragraph{Division and negation sugar.}
\begin{VNBsnip}
x / y    =>  (* x (recip y))   ; division is syntactic sugar
-x       =>  (- x)             ; unary minus; lower prec than ^
-x ^ 2   =>  (- (power x 2))  ; i.e. -(x^2), not (-x)^2
\end{VNBsnip}

\paragraph{Quantifier syntax.}
Quantifiers use the binding-list form:
\begin{VNBsnip}
forall([x], body)              ; bare variable (all classes)
forall([x in A], body)         ; guarded variable
forall([[a, b, c] in A], body) ; tuple destructuring
forall([x in A, y, z in B], body)  ; mixed bindings
forsome([x in NN], body)       ; existential; same forms
\end{VNBsnip}

\noindent
The bracket syntax after the function name distinguishes quantifiers
from ordinary function calls.

\paragraph{Display conventions (\texttt{wff->string} / \texttt{pp}).}
\begin{itemize}
\item \begin{sloppypar}\textbf{Quantifiers} print in binding-list form:
  the primitive expanded form is reverse-engineered.
  \texttt{forall([x in nn, y], p(x, y))}\end{sloppypar}
\item \textbf{Logical connectives} \texttt{and}, \texttt{or}, \texttt{iff}
  print n-ary infix with surrounding parentheses:
  \texttt{((p and q) and r)}.
  \texttt{implies} prints binary infix: \texttt{(p implies q)}.
\item \textbf{Comparisons} \texttt{in}, \texttt{=}, \texttt{==},
  \texttt{<=} print binary infix: \texttt{(x in nn)}, \texttt{(x = y)}.
\item \textbf{Arithmetic} \texttt{+}, \texttt{*} print n-ary infix:
  \texttt{(x + y + z)}.  Binary \texttt{-} prints as \texttt{(x - y)};
  unary as \texttt{-x}.  Exponentiation as \texttt{(x \textasciicircum{} y)}.
\item \textbf{Lists} \texttt{[a, b, c]} in both input and output: the
  list syntax is the same in string form (internally stored as \texttt{list}).
\item \textbf{All other compound terms} use functional notation
  \texttt{f(x, y, z)}: \texttt{nth}, \texttt{cartesian}, \texttt{fun},
  \texttt{power} (unary set power), \texttt{not}, and user-defined heads.
  Nullary symbols print without parentheses.
\end{itemize}

\paragraph{Example round-trips.}
\begin{VNBsnip}
(pp (make-wff-from-string "forall([x in NN, y in NN], x + y = y + x)"))
; => forall([x in nn, y in nn], ((x + y) = (y + x)))

(pp (make-wff-from-string "not(P) implies falsity"))
; => (not(p) implies falsity)

(pp (make-wff-from-string "2 ^ 10 = 1024"))
; => ((2 ^ 10) = 1024)

(expression->string '(forall x (implies (in x nn)
                       (forall y (implies (in y nn) (= (+ x y) (+ y x)))))))
; => "forall([x in nn, y in nn], ((x + y) = (y + x)))"
\end{VNBsnip}

\noindent
All identifiers are case-folded to lowercase by MIT~Scheme, so
\texttt{NN} and \texttt{nn} are the same symbol; \texttt{FORALL} and
\texttt{forall} are also identical.  The printer always emits lowercase.

%% =========================================================

\section{Naming the parts of a formula}
\label{sec:parts}

The parser and printer of the previous section are the surface-conforming
way to \emph{build} and \emph{display} a formula.  This section gives the
surface-conforming way to \emph{take one apart}.

A script or a piece of scratch-workspace code must never reach into the
internal S-expression by position --- \texttt{(caddr (cadr f))} and the
like.  Such code is welded to a representation accident and breaks the
moment the internal form shifts (the \texttt{and}-nesting, the
\texttt{+}/\texttt{binplus} split, a new sugar): it is the
representation-independence sin one level down.  The accessors in
\texttt{parts.scm} name the parts the way the \emph{external} syntax names
them, so the code that uses them survives changes to the internal form.

There are three entry points: \texttt{formula-kind} classifies the top
construct, \texttt{part} navigates to a named sub-part, and \texttt{match}
fills in the blanks of a surface pattern.  All three accept a wff, a
surface string, or a raw S-expression interchangeably.

\subsection{\texttt{formula-kind} --- what is this, at the top?}

\texttt{(formula-kind \textit{x})} returns a single symbol naming the
outermost surface construct, without committing to any internal detail
below it:

\begin{VNBsnip}
(formula-kind "p(x) implies q(x)")       ; => implication
(formula-kind "x = y")                   ; => equation
(formula-kind "p(x) and q(x) and r(x)")  ; => conjunction
(formula-kind "forall([x in nn], p(x))") ; => universal
(formula-kind "x in nn")                 ; => membership
(formula-kind "f(a, b)")                 ; => application
(formula-kind '((MUL R) a b))            ; => op-application
\end{VNBsnip}

\noindent
The complete set of tags follows.

\begin{itemize}
\item \texttt{truth}, \texttt{falsity} --- the logical constants.
\item \texttt{atom} --- a bare symbol; \texttt{number} --- a numeric literal.
\item \texttt{implication} (\texttt{implies}), \texttt{biconditional}
  (\texttt{iff}), \texttt{conjunction} (\texttt{and}), \texttt{disjunction}
  (\texttt{or}), \texttt{negation} (\texttt{not}) --- the connectives.
\item \texttt{equation} (\texttt{=}), \texttt{equivalence} (\texttt{==}),
  \texttt{membership} (\texttt{in}), \texttt{inequality} (\texttt{<=}),
  \texttt{subset} (\texttt{subset}) --- the binary relations.
\item \texttt{universal} (\texttt{forall}), \texttt{existential}
  (\texttt{forsome}), \texttt{description} (\texttt{iota}) --- the binders.
\item \texttt{arithmetic} (\texttt{+}, \texttt{*}).
\item \texttt{set-construction} (\texttt{union}, \texttt{intersection},
  \texttt{cartesian}).
\item \texttt{op-application} --- a structure operator such as
  \texttt{(mul R)} applied to arguments.
\item \texttt{application} --- a named predicate or function applied.
\item \texttt{unknown} --- anything that matches none of the above.
\end{itemize}

\begin{remark}
This vocabulary is meant to be \emph{closed}.  Each tag names a construct
the surface grammar actually has, so the list grows only if the grammar
itself does --- a rare event.  The pull to add a tag for some derived or
compound form should be resisted: such forms already classify as
\texttt{application}, \texttt{op-application}, or \texttt{unknown}, and a
proliferation of fine-grained tags would couple the classifier to passing
notational fashion rather than to the language.
\end{remark}

\subsection{\texttt{part} --- navigate to a sub-part}

\texttt{(part \textit{x} \textit{step} \ldots)} walks a path of selector
steps and returns the part reached.  A single result is itself a wff (it
prints through the normal verbalizer and round-trips); a \emph{plural}
selector returns a \emph{list} of wffs and must be the last step.

\begin{center}
\footnotesize
\begin{tabular}{@{}llp{6.0cm}@{}}
\toprule
Selector & Applies to & Yields \\
\midrule
\texttt{antecedent}, \texttt{consequent} & implication & the two sides \\[2pt]
\texttt{lhs}/\texttt{left}, \texttt{rhs}/\texttt{right} & any binary relation (\texttt{= == iff in <= subset}) & the two sides \\[2pt]
\texttt{negand} & negation & the negated formula \\[2pt]
\texttt{conjuncts}, \texttt{disjuncts} & conjunction / disjunction & \emph{(plural)} the flattened list of conjuncts / disjuncts \\[2pt]
\texttt{(conjunct \textit{N})}, \texttt{(disjunct \textit{N})} & conjunction / disjunction & the \textit{N}-th, 1-indexed \\[2pt]
\texttt{binder} & quantifier (\texttt{forall}/\texttt{forsome}/\texttt{iota}) & the bound variable \\[2pt]
\texttt{matrix}/\texttt{scope}/\texttt{body} & quantifier & the scope \\[2pt]
\texttt{operator}/\texttt{op}/\texttt{head} & application & the operator (e.g.\ \texttt{(mul r)}) \\[2pt]
\texttt{operands}/\texttt{args} & application & \emph{(plural)} the argument list \\[2pt]
\texttt{(operand \textit{N})} & application & the \textit{N}-th argument, 1-indexed \\
\bottomrule
\end{tabular}
\end{center}

\begin{VNBsnip}
(part "p(x) implies q(x)" 'antecedent)         ; => p(x)
(part "x = y" 'lhs)                            ; => x
(part "p(x) and q(x) and r(x)" '(conjunct 2))  ; => q(x)
(part "p(x) and q(x) and r(x)" 'conjuncts)     ; => (p(x) q(x) r(x))
(part "f(a, b, c)" '(operand 2))               ; => b
(part "f(a, b, c)" 'operands)                  ; => (a b c)
(part "forall([x in nn], p(x))" 'binder)       ; => x
(part "forall([x in nn], p(x))" 'matrix)       ; => (x in nn implies p(x))

;; steps chain left-to-right: descend, then descend again
(part "p(x) implies y = z" 'consequent 'lhs)   ; => y
\end{VNBsnip}

\noindent
The selectors are \emph{total over their declared shape and loud
otherwise}: \texttt{(part "x = y" 'antecedent)} is an \emph{error}, not a
silent \texttt{\#f} --- the same partial-\texttt{=} definedness discipline
that runs through the rest of the system (Section~\ref{sec:funapp}).

\begin{remark}
The accessors honour the \emph{surface structure} of a formula --- the
string-syntax view laid out in Section~\ref{sec:parser-printer} --- not the
internal S-expression that stores it.  \texttt{conjuncts} flattens the
associative spine, so \texttt{p and q and r} yields three conjuncts however
the binary \texttt{and} happens to be nested.  \texttt{operands} de-tuples the apply-tupling convention
(Section~\ref{sec:curried-tupled-apply}), so the single internal
\texttt{(list a b)} argument comes back as the operand list
\texttt{(a~b)}.  And the binary structure-op application \texttt{(mul
R)(a, b)} is first-class: \texttt{operator} gives \texttt{(mul r)},
\texttt{operands} gives \texttt{(a b)}.
\end{remark}

\subsection{\texttt{match} --- fill the blanks of a pattern}

\texttt{(match \textit{pattern-string} \textit{x})} matches a surface
pattern against \textit{x}.  The pattern is an ordinary formula in which
some sub-parts are left \emph{blank}, each blank written as a name with a
leading question mark --- \texttt{?P}, \texttt{?Q}, \texttt{?a} --- standing
for an arbitrary sub-formula or term.  The call returns an association list
\texttt{((?v\ .\ \textit{wff})\ \ldots)} pairing each blank with what it
matched, or \texttt{\#f} if the pattern does not fit.

\begin{VNBsnip}
(match "?P implies ?Q" "x in nn implies x in zz")
; => ((?q . x in zz) (?p . x in nn))

(match "?P and ?Q" "x in nn implies x in zz")
; => #f
\end{VNBsnip}

\noindent
Matching reuses the macete one-way matcher
(Section~\ref{sec:macetes}) rather than introducing a second pattern
engine, and binder desugaring is applied to both sides so a quantified
pattern matches a quantified target.  Where \texttt{part} is the right tool
when the shape is known in advance, \texttt{match} is the right tool when
the shape is the thing being tested.

\begin{remark}
Input is coerced uniformly: \texttt{part} and friends accept a wff, a
surface string, or a raw S-expression.  Given a wff or string they hand
back wffs (re-wrapped in the same library); given a bare S-expression there
is no wff parent to re-wrap against, so the parts come back as raw
S-expressions.  One asymmetry follows from \texttt{make-wff} validation: a
standalone \emph{term} string such as \texttt{"x + y"} is not a wff, so
\texttt{(formula-kind "x + y")} signals an error --- classify a term by
passing the S-expression, \texttt{(formula-kind '(+ x y))}, or reach it as
a \texttt{part} of the formula it sits in.
\end{remark}

%% =========================================================


%% =========================================================
\section{Formulas as typed objects: \texttt{wff}}
\label{chap:syntax-contexts}

A \emph{well-formed formula} in this system is not a bare S-expression
but a typed object: a formula that has passed the well-formedness check.  In the REPL it prints as:
\begin{VNBsnip}
#[wff N hD "formula string"]
\end{VNBsnip}
where \texttt{N} is MIT Scheme's object hash, \texttt{hD} is the
formula's digest in hexadecimal, and the quoted string is the
pretty-printed formula.  Every formula in a sequent is a \texttt{wff}.

\texttt{(make-wff expr)} constructs such an object, from a string in
surface syntax or from an S-expression.  \texttt{(free-vars (wff->sexp p))}
returns the variables free in the formula: \texttt{free-vars} works on
the expression, not on the \texttt{wff} record.

Two procedures expose the internals of a \texttt{wff}:
\begin{VNBsnip}
(wff->sexp p)        ; => raw S-expression (no output)
(display-contents p) ; prints kind, library name, formula
\end{VNBsnip}
\texttt{wff->sexp} is defined in \texttt{wff.scm};
\texttt{display-contents} is defined in \texttt{contexts.scm}.
Both are available in interactive sessions.

Example (the base library is \texttt{VNB-LIBRARY}):
\begin{VNBsnip}
(define p (make-wff "forall([x in A], 0 <= funny-name)"))
; => #[wff 12 "forall([x in a], 0 <= funny-name)"]

(free-vars (wff->sexp p))
; => (a funny-name)   ; x is bound by the forall
\end{VNBsnip}


%% =========================================================
\section{Notation: \texttt{notation!}}
\label{sec:notation}

Every head the system knows --- a predicate, a functoid, a structure predicate,
even a primitive like \texttt{=} or \texttt{in} --- has an entry in a single
operator table.  \texttt{def-predicate}, \texttt{def-functoid} and
\texttt{declare-structure} fill in the \emph{structural} part of the entry (its
kind, its arity) at the definition point.  \texttt{notation!} fills in the
\emph{reading} --- how the head renders in \LaTeX, in English, and to the
parser.  The two are orthogonal: re-declaring a head updates the structural slots
and keeps whatever notation was declared, so a \texttt{notation!} may sit beside
the definition or well after it.

It is a keyword form; a call names only the slots it has to set:
\begin{VNBcode}
(notation! 'IS-RING 'noun "ring" 'article "a")
(notation! 'IS-EUCLIDEAN-RING 'kind 'predicate 'arity 1
           'noun "Euclidean ring" 'article "a")
(notation! 'IS-COUNTABLE-PSEUDOMETRIC-FAMILY 'kind 'predicate 'arity 2
           'english "$1 is a countable family of pseudometrics on $2")
\end{VNBcode}

The recognised keys are
\texttt{tex}, \texttt{english}, \texttt{noun}, \texttt{article},
\texttt{prec}, \texttt{infix?}, \texttt{arity}, and \texttt{kind}.

\begin{itemize}
\item \texttt{tex} and \texttt{english} are \emph{templates} --- strings with
  placeholders \texttt{\$1}\ldots\texttt{\$9} for the rendered arguments.  A
  binary operator combining both with \texttt{prec} (a precedence for
  parenthesisation) and \texttt{infix?} reads, schematically,
  \texttt{(notation! '+ 'tex "\$1 + \$2" 'english "\$1 plus \$2" 'prec 60 'infix?
  \#t)}.  When a template cannot express the rendering, \texttt{english} may
  instead be a \emph{procedure} from the list of argument strings to a string.
\item \texttt{noun} and \texttt{article} are the special case of a unary
  structure predicate: given them, \texttt{is-ring($s$)} reads ``$s$ is a
  ring'' throughout the prose renderers, and the article joins correctly
  (``a ring'', ``an $R$-module'').
\item \texttt{kind} and \texttt{arity} let a \texttt{notation!} register a head
  the def-forms never introduced --- a primitive such as \texttt{=} or
  \texttt{in} --- so it, too, gets a reading.
\end{itemize}

One table, one declaration form, read by every renderer: \texttt{expr->tex} and
\texttt{wff-english} for formulas, \texttt{describe-structure} for the structure
cards, and the generated \texttt{OPERATORS.md}.  A head with no
\texttt{notation!} still works --- it falls back to prefix application,
\texttt{is-ring($s$)} --- so notation is an addition, never a prerequisite.

%%% -----------------------------------------------------------------------
\section{The optional infix layer}
\label{sec:infix-layer}

The TeX printer writes an applied structure operation in functional
form: \texttt{((add K) a b)} prints as $\mathtt{add}(K)(a, b)$.  The
\emph{infix layer} (\texttt{tex-infix.scm}) writes it instead as
\[
  a +_K b .
\]
It is off by default, and while it is off the printer writes the
functional form.  It changes only TeX output: the panels, the
proof reader and \texttt{display-proof}.  The parser, \texttt{pp} and
the kernel never see it.

It differs from \texttt{notation!} (Section~\ref{sec:notation}) in
two ways.  A \texttt{notation!} declares the reading of one head and is
always in force.  The layer reads an \emph{accessor applied to a
structure}, so one row serves \texttt{add} in every ring, module and
field, and it is switched on and off as a whole.

\subsection{Turning it on}

In the workspace, the menu \textsf{VNB > Printer Macros} has the
toggle \textsf{Printer Macros On}; the panels repaint at once.  At the
REPL:
\begin{VNBcode}
(tex-infix! #t)                          ; on for every printout
(tex-infix! #f)                          ; off
(display-proof 'NAME "file.pdf" 'infix)  ; on for this printout only
(infix-tex 'NAME)                        ; the TeX of a theorem, layer on
(with-tex-infix CONV thunk)              ; one printout, conventions CONV
\end{VNBcode}

\subsection{What it prints}

An accessor applied to a structure term $S$, then to its arguments,
is written with the subscript $S$:
\[
\begin{array}{ll@{\qquad}ll}
\texttt{((add S) a b)}  & a +_S b        & \texttt{((mul S) a b)}  & a \times_S b \\
\texttt{((opr S) a b)}  & a \cdot_S b    & \texttt{((act S) c v)}  & c \cdot_S v \\
\texttt{((neg S) a)}    & -_S\, a        & \texttt{((inv S) a)}    & a^{-1}_S \\
\texttt{(zero S)}       & 0_S            & \texttt{(iden S)}       & e_S \\
\texttt{((fnrm S) a)}   & \lvert a \rvert_S & \texttt{((vnrm S) v)} & \lVert v \rVert_S \\
\texttt{((ip S) u v)}   & \langle u, v \rangle_S & \texttt{((vstar S) x)} & x^{*}_S
\end{array}
\]
The full list is the base table (below).  $S$ is read through views:
a \texttt{def-functor} view such as
\texttt{(normed-field-as-commutative-ring K)} is written $K$, since the
view's slots are $K$'s slots.

A defined head can also have a template.  The shipped templates print
the sums and products over an index:
\[
  \texttt{series-partial-sum(c, n)} \quad\text{as}\quad \sum_{i < n} c(i).
\]
When the summand is a \texttt{vnb-lambda} over $\mathtt{nn}$, the
lambda's own variable is used and its body printed, as in
$\sum_{k < n} c(k)\, x^k$.  Otherwise the printer invents an index
variable that clashes with no symbol of the formula.

\subsection{Faithfulness}

The layer changes notation, never the term.  Nothing is simplified or
rewritten: $a + (-b)$ prints as
\[
  a +_S -_S\, b ,
\]
never as $a - b$.  Every subterm the plain printer would show is still
there.

Parentheses follow the binding strengths: $+$ is weakest, then the
products, then unary minus, then the postfix operations.  A chain of
one associative operation of one structure nested to the \emph{left}
drops its parentheses.  Nested to the right it keeps them, so an
associativity law still reads as one:
\[
  a + b + c \;=\; a + (b + c).
\]

\subsection{Conventions: when the subscript is dropped}

The conventions are a list, like a typing block:
\begin{itemize}
\item \texttt{(auto)}, the default.  Drop the subscript when one
  structure is in play in the formula being printed and it is a
  variable, not a named constant such as \texttt{zz-ring}.  Its
  constants $0_S$, $1_S$, $e_S$ keep their subscript when the formula
  also holds a numeral, so $0_K$ never reads as the number $0$.
\item \texttt{(keep)}.  Always write the subscript.
\item \texttt{(drop S \ldots)}.  Drop the subscript of these structure
  terms.
\item \texttt{(symbol HEAD "tex")}.  Write \texttt{HEAD}'s operation
  with this symbol.
\end{itemize}
``The formula being printed'' is the term given to the outermost
printer call.  A display broken over rows is scanned row by row.

\subsection{The tables}

The operations, their TeX and the default conventions are data, not
code.  The base table is \texttt{notation/base.tbl} in the tree; its
head describes the format.  A row reads
\begin{VNBcode}
; KIND     HEAD      PREC  ASSOC  TEX
infix      add       10    assoc  +
postfix    inv       40    -      ^{-1}
outfix     vnrm      100   -      \norm{#1}
apply      series-partial-sum  5  -  \sum_{#i < #2} #1(#i)
\end{VNBcode}
The kinds are \texttt{infix}, \texttt{prefix}, \texttt{postfix},
\texttt{const}, \texttt{outfix} (a template in \texttt{\#1}),
\texttt{bracket} (in \texttt{\#1} and \texttt{\#2}) and \texttt{apply}
(a template for a defined head).  An \texttt{apply} template must use
every argument: it may not hide one.  There are also \texttt{macro}
rows, \texttt{off HEAD}, and \texttt{convention} rows.

Further tables are read after the base table, and a later row replaces
an earlier row of the same head.  The order is:
\begin{enumerate}
\item the tables named, one per line, in
  \texttt{\textasciitilde/vnb-notation/active}.  A name is
  \texttt{\textasciitilde/vnb-notation/NAME.tbl} when that file exists,
  else the tree's \texttt{notation/NAME.tbl};
\item the personal table \texttt{\textasciitilde/vnb-notation/personal.tbl},
  when it exists.
\end{enumerate}
The menu \textsf{VNB > Printer Macros} turns each table on or off, and
\textsf{Edit or New Table\ldots} and \textsf{Edit Personal Table} open
one for editing.

A changed table is re-read at the next printout with the layer on, or
at once by \textsf{Reload} or \texttt{(tex-infix-reload!)}.  A reload
with any row in error changes nothing.  The errors are listed by file
and line, and the table in force stays.
